\documentclass[11pt,a4paper]{article}
\usepackage[utf8]{inputenc}
\usepackage[T1]{fontenc}
\usepackage{amsmath,amsfonts,amssymb}
\usepackage{graphicx}
\usepackage{hyperref}
\usepackage{listings}
\usepackage{xcolor}
\usepackage{booktabs}
\usepackage{float}
\usepackage{geometry}
\geometry{margin=1in}

% Color definitions
\definecolor{zenblue}{RGB}{41,121,255}
\definecolor{zengreen}{RGB}{52,199,89}
\definecolor{zenorange}{RGB}{255,149,0}
\definecolor{codegray}{RGB}{245,245,245}

% Hyperref setup
\hypersetup{
    colorlinks=true,
    linkcolor=zenblue,
    urlcolor=zenblue,
    citecolor=zenblue
}

% Code listing setup
\lstset{
    backgroundcolor=\color{codegray},
    basicstyle=\ttfamily\small,
    breaklines=true,
    captionpos=b,
    frame=single,
    numbers=left,
    numberstyle=\tiny\color{gray}
}

\title{
    \vspace{-2cm}
    \Large \textbf{Zen AI Model Family} \\
    \vspace{0.5cm}
    \Huge \textbf{Zen-Designer-235B-A22B-Thinking} \\
    \vspace{0.3cm}
    \large Vision-Language Reasoning for Design \\
    \vspace{0.5cm}
    \normalsize Technical Whitepaper v1.0
}

\author{
    Zach Kelling\thanks{zach@lux.network} \\
    \texttt{research@hanzo.ai} \\
    \\
    Zoo Labs Foundation \\
    \texttt{foundation@zoolabs.org}
}

\date{September 2025}

\begin{document}

\maketitle

\begin{abstract}
We present \textbf{Zen-Designer-235B-A22B-Thinking}, a deliberative vision-language model for
design reasoning and visual analysis. The model is a derivative of the open-weight
\textbf{Qwen3-VL-235B-A22B} mixture-of-experts foundation (Apache~2.0), adapted by Hanzo AI and
the Zoo Labs Foundation to emit explicit intermediate reasoning before producing a final answer
on visual and layout tasks. It is a sparse MoE with 235B total parameters of which approximately
22B are active per token, served under the \texttt{Qwen3VLForConditionalGeneration} architecture.
Weights are released openly under Apache~2.0 at
\href{https://huggingface.co/zenlm/zen-designer-235b-a22b-thinking}{huggingface.co/zenlm/zen-designer-235b-a22b-thinking}.
This whitepaper documents the architecture as published in the released configuration and the
intended use of the thinking variant; it does not claim novel pretraining or report benchmark
numbers that have not been independently reproduced.
\end{abstract}

\tableofcontents
\newpage

\section{Introduction}

\textbf{Zen-Designer-235B-A22B-Thinking} is the deliberative member of the Zen-Designer line. It
is tuned to produce a chain of intermediate reasoning steps before a final response, trading
latency for improved performance on multi-step visual-reasoning tasks. It is built on the
open-weight \textbf{Qwen3-VL-235B-A22B} foundation released by the Qwen team under the Apache~2.0
license. A companion \emph{instruct} variant
(\href{https://huggingface.co/zenlm/zen-designer-235b-a22b-instruct}{zen-designer-235b-a22b-instruct})
provides lower-latency, direct responses.

\subsection{Provenance and Licensing}
This model is a derivative of an upstream open-weight foundation. We state the lineage
explicitly to avoid any overclaiming:
\begin{itemize}
    \item \textbf{Upstream foundation}: Qwen3-VL-235B-A22B (vision-language MoE), Apache~2.0.
    \item \textbf{Derivative}: Zen-Designer-235B-A22B-Thinking, reasoning-tuned for design tasks.
    \item \textbf{License}: Apache~2.0 (inherited; redistribution permitted with attribution).
    \item \textbf{Repository}: \texttt{zenlm/zen-designer-235b-a22b-thinking} on Hugging Face.
\end{itemize}

\subsection{Scope}
\begin{itemize}
    \item \textbf{Sparse MoE efficiency}: $\sim$22B parameters active per token from a 235B total pool.
    \item \textbf{Deliberative reasoning}: explicit intermediate steps before a final answer.
    \item \textbf{Vision-language}: image-conditioned reasoning and text-conditioned design assistance.
    \item \textbf{Extended context}: 131K-token maximum position embeddings (per released config).
\end{itemize}

\section{Architecture}

\subsection{Model Configuration}

The values below are taken directly from the released model configuration
(\texttt{config.json}); they describe the published architecture rather than any internally
measured performance.

\begin{table}[H]
\centering
\begin{tabular}{ll}
\toprule
\textbf{Component} & \textbf{Specification} \\
\midrule
Architecture & \texttt{Qwen3VLForConditionalGeneration} \\
Total Parameters & 235B (sparse MoE) \\
Active Parameters & $\sim$22B per token \\
Experts / Active per Token & 64 / 4 \\
Text Hidden Size & 8192 \\
Text Layers & 80 \\
Attention Heads / KV Heads & 64 / 8 \\
Vision Hidden Size & 2048 \\
Vision Layers & 48 \\
Vision Patch / Image Size & 14 / 2048 \\
Vocabulary Size & 151{,}936 \\
Max Position Embeddings & 131{,}072 \\
Precision & bfloat16 \\
\bottomrule
\end{tabular}
\caption{Zen-Designer-235B-A22B-Thinking configuration (from released \texttt{config.json}).}
\end{table}

\subsection{Mixture of Experts}
The text backbone is a sparsely activated mixture-of-experts transformer: each token is routed
to a small subset (top-4 of 64) of expert feed-forward blocks per layer, so that only a fraction
of the 235B total parameters is exercised on any given forward pass. This is the standard MoE
efficiency trade-off described by Shazeer et al.~\cite{shazeer2017outrageously} and
Fedus et al.~\cite{fedus2022switch}, inherited from the Qwen3-VL foundation.

\subsection{Deliberative Reasoning}
The thinking variant is tuned to emit an explicit reasoning trace prior to its final response,
following the chain-of-thought paradigm~\cite{wei2022chain}. This typically improves accuracy on
multi-step visual and layout-reasoning tasks at the cost of additional generated tokens and
latency relative to the instruct variant.

\subsection{Vision-Language Integration}
A vision encoder (48 layers, hidden size 2048, patch size 14) produces visual tokens that are
consumed by the text decoder under the \texttt{Qwen3VLForConditionalGeneration} interface,
enabling image-conditioned reasoning alongside text-only generation.

\section{Intended Use}

\subsection{Primary Applications}
\begin{itemize}
    \item Multi-step UI/UX design analysis and critique
    \item Layout and information-architecture reasoning
    \item Visual question answering requiring step-by-step justification
    \item Design-system review with explicit rationale
    \item Accessibility evaluation with documented reasoning
\end{itemize}

\subsection{Thinking vs.\ Instruct}
The thinking variant is preferred where multi-step visual reasoning or self-verification adds
value and additional latency is acceptable. For low-latency, direct responses, use the
companion instruct variant.

\subsection{Usage Example}

\begin{lstlisting}[language=Python, caption=Loading and inference]
from transformers import AutoModelForVision2Seq, AutoProcessor

model = AutoModelForVision2Seq.from_pretrained(
    "zenlm/zen-designer-235b-a22b-thinking")
processor = AutoProcessor.from_pretrained(
    "zenlm/zen-designer-235b-a22b-thinking")

# Image-conditioned reasoning
inputs = processor(images=image, text="Analyze this UI", return_tensors="pt")
outputs = model.generate(**inputs)
analysis = processor.decode(outputs[0])
\end{lstlisting}

\section{Deployment}

\subsection{Available Formats}
\begin{itemize}
    \item \textbf{SafeTensors}: native bfloat16 weights (sharded).
    \item \textbf{GGUF / MLX}: community quantizations may be provided separately.
\end{itemize}

\subsection{Hardware Footprint}
Because the model carries 235B total parameters, full-precision serving requires
multi-accelerator deployment; quantization reduces this footprint at some quality cost. The
thinking variant additionally emits more tokens per query than the instruct variant, which
increases end-to-end latency. We do not publish a single canonical latency figure, as
throughput depends heavily on the serving stack, batch size, quantization, and hardware.

\section{Limitations}

\begin{itemize}
    \item \textbf{Inherited behavior}: as a derivative of Qwen3-VL, the model inherits the
          capabilities, biases, and limitations of its upstream foundation and training data.
    \item \textbf{No independent benchmark claims}: this document intentionally omits benchmark
          scores that have not been independently reproduced for this derivative.
    \item \textbf{Latency}: the deliberative reasoning trace increases token count and latency
          relative to the instruct variant.
    \item \textbf{Resource intensity}: the 235B parameter count makes self-hosting costly
          relative to smaller members of the Zen family.
    \item \textbf{Visual generation}: the model assists with design \emph{reasoning} and
          analysis; pixel-level image synthesis is out of scope for the language backbone.
\end{itemize}

\section{Conclusion}

\textbf{Zen-Designer-235B-A22B-Thinking} packages an open-weight Qwen3-VL-235B-A22B foundation as
a deliberative design and visual-reasoning assistant, released openly under Apache~2.0. We have
described its architecture as published and stated its provenance plainly so that downstream
users can make informed decisions about licensing, capability, and cost.

\section*{Acknowledgments}

We thank the Qwen team for releasing the Qwen3-VL foundation under Apache~2.0, and the
open-source community and the teams at Hanzo AI and Zoo Labs Foundation for their contributions
to this work.

\begin{thebibliography}{99}
\bibitem{vaswani2017attention} Vaswani, A. et al. (2017). Attention Is All You Need. NeurIPS 2017.
\bibitem{radford2021learning} Radford, A. et al. (2021). Learning Transferable Visual Models From Natural Language Supervision. ICML 2021.
\bibitem{shazeer2017outrageously} Shazeer, N. et al. (2017). Outrageously Large Neural Networks: The Sparsely-Gated Mixture-of-Experts Layer. ICLR 2017.
\bibitem{fedus2022switch} Fedus, W. et al. (2022). Switch Transformers: Scaling to Trillion Parameter Models with Simple and Efficient Sparsity. JMLR 2022.
\bibitem{wei2022chain} Wei, J. et al. (2022). Chain-of-Thought Prompting Elicits Reasoning in Large Language Models. NeurIPS 2022.
\bibitem{qwen3vl2025} Qwen Team (2025). Qwen3-VL: Vision-Language Models. \href{https://huggingface.co/Qwen}{huggingface.co/Qwen}.
\end{thebibliography}

\appendix

\section{Model Card}

\begin{table}[H]
\centering
\begin{tabular}{ll}
\toprule
\textbf{Field} & \textbf{Value} \\
\midrule
Model Name & Zen-Designer-235B-A22B-Thinking \\
Upstream Foundation & Qwen3-VL-235B-A22B (Apache 2.0) \\
Version & 1.0.0 \\
Release Date & September 2025 \\
License & Apache 2.0 \\
Repository & \href{https://huggingface.co/zenlm/zen-designer-235b-a22b-thinking}{huggingface.co/zenlm/zen-designer-235b-a22b-thinking} \\
Documentation & \href{https://github.com/zenlm/zen}{github.com/zenlm/zen} \\
Contact & research@hanzo.ai \\
\bottomrule
\end{tabular}
\caption{Model Card Information}
\end{table}

\end{document}
